%Data institusi

%\input{Dozentdaten}

%\renewcommand{\fachbereich}{Fachbereich}

%\renewcommand{\dozent}{Prof. Dr. . }

%Data ujian

\renewcommand{\klausurgebiet}{ }

\renewcommand{\klausurtyp}{ }

\renewcommand{\klausurdatum}{ . 20}

\klausurvorspann {\fachbereich} {\klausurdatum} {\dozent} {\klausurgebiet} {\klausurtyp}

%Data untuk tabel nilai berikut

\renewcommand{\aeins}{  3  }

\renewcommand{\azwei}{  3   }

\renewcommand{\adrei}{  8  }

\renewcommand{\avier}{  7  }

\renewcommand{\afuenf}{  8  }

\renewcommand{\asechs}{  3  }

\renewcommand{\asieben}{  5  }

\renewcommand{\aacht}{  5  }

\renewcommand{\aneun}{  6  }

\renewcommand{\azehn}{  4  }

\renewcommand{\aelf}{  12  }

\renewcommand{\azwoelf}{  64  }

\renewcommand{\adreizehn}{     }

\renewcommand{\avierzehn}{    }

\renewcommand{\afuenfzehn}{    }

\renewcommand{\asechzehn}{    }

\renewcommand{\asiebzehn}{    }

\renewcommand{\aachtzehn}{    }

\renewcommand{\aneunzehn}{    }

\renewcommand{\azwanzig}{    }

\renewcommand{\aeinundzwanzig}{    }

\renewcommand{\azweiundzwanzig}{    }

\renewcommand{\adreiundzwanzig}{    }

\renewcommand{\avierundzwanzig}{    }

\renewcommand{\afuenfundzwanzig}{    }

\renewcommand{\asechsundzwanzig}{    }

\punktetabelleelf

\klausurnote

\newpage

\setcounter{section}{0}

\inputaufgabegibtloesung
{3}
{

Definisikan istilah-istilah berikut
\zusatzklammer {yang dicetak miring} {} {.}
\aufzaehlungsechs{Suatu
\stichwort {medan normal} {}
$N$ pada hiperpermukaan diferensiabel

\mavergleichskette
{\vergleichskette
{ Y
}
{ \subseteq }{ \R^n
}
{  }{
}
{    }{
}
{   }{
}
}
{}{}{.}

}{Suatu ruang topologis yang
\stichwort {parakompak}{}
.

}{
\stichwort {Ruang kotangen} {}
di suatu titik

\mavergleichskette
{\vergleichskette
{  P
}
{ \in }{ M
}
{  }{}
{    }{}
{   }{}
}
{}{}{}
dari
\definitionsverweis {manifold diferensiabel}{}{}
$M$.

}{Suatu
\stichwort {titik reguler} {}

\mathl{Q \in L}{} bagi
\definitionsverweis {pemetaan diferensiabel}{}{}
\maabbdisp {\varphi} { L } { M
} {}
antara
\definitionsverweis {manifold-manifold diferensiabel}{}{}
\mathkor {} {L} {dan} {M} {.}

}{
\stichwort {Produk} {}
dari manifold-manifold diferensiabel
\mathkor {} {M} {dan} {N} {.}

}{Suatu
\stichwort {koneksi} {}
pada bundel vektor diferensiabel
\maabb {} {E} {M
} {}
di atas manifold diferensiabel $M$.
}

}
{} {}

\inputaufgabegibtloesung
{3}
{

Nyatakan teorema-teorema berikut.
\aufzaehlungdrei{
\stichwort {Teorema tentang hubungan antara kelengkungan dan vektor normal satuan suatu kurva bidang berparameter panjang busur} {.}}{
\stichwort {Theorema egregium} {.}}{
\stichwort {Teorema Green} {}
untuk luas.}

}
{} {}

\inputaufgabegibtloesung
{8}
{

Misalkan

\mavergleichskettedisp
{\vergleichskette
{ f(x,y)
}
{ =} { ax^2+by^2+cxy
}
{ } {
}
{ } {
}
{ } {
}
}
{}{}{.}
Tentukan
\definitionsverweis {kelengkungan-kelengkungan utama}{}{}
dan
\definitionsverweis {arah-arah kelengkungan utama}{}{}
dari
\definitionsverweis {grafik}{}{}
$f$ di titik asal.

}
{} {}

\inputaufgabegibtloesung
{7}
{

Buktikan teorema eksistensi untuk transpor paralel pada suatu hiperpermukaan

\mavergleichskette
{\vergleichskette
{ Y
}
{ \subseteq }{ \R^n
}
{  }{
}
{    }{
}
{   }{
}
}
{}{}{.}

}
{} {}

\inputaufgabegibtloesung
{8}
{

Tunjukkan bahwa suatu
\definitionsverweis {manifold diferensiabel}{}{}
$M$ berdimensi

\mavergleichskette
{\vergleichskette
{ n
}
{ \geq }{ 1
}
{  }{
}
{    }{
}
{   }{
}
}
{}{}{}
mempunyai tak hingga banyak
\definitionsverweis {difeomorfisme}{}{}
\maabbdisp {\varphi} { M } { M
} {}
.

}
{} {}

\inputaufgabegibtloesung
{3}
{

Misalkan $M$ suatu
\definitionsverweis {manifold diferensiabel}{}{}
dan
\maabbdisp {f} { M } {\R
} {}
suatu
\definitionsverweis {fungsi terdiferensialkan secara kontinu}{}{.}
Misalkan di titik

\mavergleichskette
{\vergleichskette
{ P
}
{ \in }{ M
}
{  }{
}
{    }{
}
{   }{
}
}
{}{}{}
fungsi tersebut mempunyai
\definitionsverweis {ekstremum lokal}{}{.}
Tunjukkan bahwa

\mavergleichskettedisp
{\vergleichskette
{ (df)_P
}
{ =} { 0
}
{ } {
}
{ } {
}
{ } {
}
}
{}{}{.}

}
{} {}

\inputaufgabegibtloesung
{5}
{

Kita tinjau
$1$-\definitionsverweis {bentuk diferensial}{}{}

\mavergleichskettedisp
{\vergleichskette
{ \omega
}
{ =} { -vdu+udv
}
{ } {
}
{ } {
}
{ } {
}
}
{}{}{}
pada
$1$-\definitionsverweis {bola}{}{}

\mavergleichskette
{\vergleichskette
{ S^1
}
{ \subset }{ \R^2
}
{  }{
}
{    }{
}
{   }{
}
}
{}{}{,}
dengan koordinat ruang sekitar $\R^2$ dinotasikan oleh
\mathkor {} {u} {dan} {v} {.}
Tentukan
\mathl{\pi^* \omega}{} di bawah
\definitionsverweis {pemetaan terdiferensialkan secara kontinu}{}{}
\maabbeledisp {\pi} {\R^2 \setminus \{0\}} { S^{1}
} {(x , y) } { { \frac{   1  }{  \sqrt{x^2 + y^2}  } } (x , y)
} {.}

}
{} {}

\inputaufgabegibtloesung
{5}
{

Buktikan teorema tentang perhitungan volume kanonik pada manifold Riemann $M$.

}
{} {}

\inputaufgabegibtloesung
{6 (1+1+3+1)}
{

Jelaskan mengapa $\R^2$ tidak membawa struktur
\definitionsverweis {manifold berbatas}{}{}
sedemikian sehingga subhimpunan $T$ yang diberikan merupakan
\definitionsverweis {batas}{}{.}
\aufzaehlungvierabc{
\mavergleichskette
{\vergleichskette
{ T
}
{ = }{ {]0,1[}
}
{  }{
}
{    }{
}
{   }{
}
}
{}{}{,}
dengan interval tersebut terletak pada sumbu $x$.
}{
\mavergleichskette
{\vergleichskette
{ T
}
{ = }{ [0,1]
}
{  }{
}
{    }{
}
{   }{
}
}
{}{}{,}
dengan interval tersebut terletak pada sumbu $x$.
}{
\mavergleichskette
{\vergleichskette
{ T
}
{ = }{ \R
}
{  }{
}
{    }{
}
{   }{
}
}
{}{}{,}
dengan $\R$ adalah sumbu $x$.
}{
\mavergleichskette
{\vergleichskette
{ T
}
{ = }{ S^1
}
{  }{
}
{    }{
}
{   }{
}
}
{}{}{.}
}

}
{} {}

\inputaufgabegibtloesung
{4}
{

Misalkan $\omega$ suatu
$n-1$-\definitionsverweis {bentuk diferensial}{}{}
yang terdiferensialkan secara kontinu pada
\definitionsverweis {himpunan terbuka}{}{}

\mavergleichskette
{\vergleichskette
{ U
}
{ \subseteq }{ \R^n
}
{  }{
}
{    }{
}
{   }{
}
}
{}{}{}
dan misalkan
\definitionsverweis {turunan eksterior}{}{}
sama dengan

\mavergleichskettedisp
{\vergleichskette
{d \omega
}
{ =} { fdx_1 \wedge \ldots \wedge dx_n
}
{ } {
}
{ } {
}
{ } {
}
}
{}{}{}
dengan suatu fungsi
\maabb {f} { U } {\R
} {.}
Tunjukkan untuk

\mavergleichskette
{\vergleichskette
{ P
}
{ \in }{ U
}
{  }{
}
{    }{
}
{   }{
}
}
{}{}{}
kesamaan

\mavergleichskettedisp
{\vergleichskette
{ f(P)
}
{ =} { { \frac{   1  }{  \beta_n  } } \cdot \operatorname{lim}_{   \epsilon \rightarrow  0  } \, { \frac{   1  }{  \epsilon^n } } \int_{S^{n-1}(P, \epsilon)} \omega
}
{ } {
}
{ } {
}
{ } {
}
}
{}{}{,}
dengan $\beta_n$ menyatakan
volume
bola satuan berdimensi $n$.

}
{} {}

\inputaufgabegibtloesung
{12}
{

Buktikan teorema Stokes untuk manifold berbatas.

}
{} {}
